% -------------------------------------------------------------------------
% WBS ID: RES-DAT-UNC
% Deliverable: Data Quality, Uncertainty and Provenance
% Status: In progress
% Evidence status: Data provenance and preprocessing uncertainty requirements established
% Review status: Not reviewed
% Last updated: 2026-07-19
% Dependencies: RES-DAT-GEO, RES-DAT-SRC, RES-DAT-OBS, RES-DAT-SCN
% Downstream interfaces: RES-DAT-CAS, RES-VER-MET, Software data ingestion
% -------------------------------------------------------------------------

\section{RES-DAT-UNC --- Data Quality, Uncertainty and Provenance}
\label{sec:res-dat-unc}

\subsection{Purpose and Scope}
\label{sec:res-dat-unc-purpose}

This Deliverable defines the provenance, CRS, datum, resolution,
licence and preprocessing uncertainty records required for all Tōhoku
case-study data. It covers uncertainty traceability for source,
topobathymetry, observations, corridor geometry and derived replay
boundaries. It does not quantify final validation tolerances; those
belong to `RES-VER-MET`.

\subsection{Research Questions}
\label{sec:res-dat-unc-research-questions}

The uncertainty decision asks:

\begin{enumerate}
  \item Which metadata are mandatory before a dataset can be ingested?
  \item Which preprocessing operations change the meaning or error
    structure of the data?
  \item Which uncertainty flags must travel into validation reports and
    local replay boundaries?
\end{enumerate}

\subsection{Terminology and Definitions}
\label{sec:res-dat-unc-definitions}

% TODO: Define the technical terminology, symbols, quantities and
% classifications used in this Deliverable.

\subsection{Literature Evidence}
\label{sec:res-dat-unc-literature-evidence}

% TODO: Record evidence from peer-reviewed papers, authoritative datasets,
% standards, technical reports and primary documentation.
%
% Do not create a detached literature-summary list. Explain how each source
% supports, contradicts or limits a project decision.

\subsection{Governing Relations and Quantitative Definitions}
\label{sec:res-dat-unc-governing-relations}

% TODO: Record relevant equations, dimensionless groups, empirical models,
% extraction rules or quantitative definitions.
%
% Use align environments for displayed equations.
% Every displayed equation must have a unique \label{eq:...}.
% Do not place punctuation after displayed equations.

Each data asset $d$ shall carry a provenance record:

\begin{align}
  \mathcal{P}_{d}
  &=
  \left\{
    s_d,\,
    c_d,\,
    \mathcal{R}_{h,d},\,
    \mathcal{R}_{v,d},\,
    \Delta_d,\,
    \ell_d,\,
    \mathcal{T}_d,\,
    \mathcal{U}_d
  \right\}
  \label{eq:res-dat-unc-provenance-record}
\end{align}

where $s_d$ is source/provider, $c_d$ is citation key,
$\mathcal{R}_{h,d}$ is horizontal CRS, $\mathcal{R}_{v,d}$ is vertical
datum, $\Delta_d$ is native resolution, $\ell_d$ is licence/usage note,
$\mathcal{T}_d$ is the ordered preprocessing transform and
$\mathcal{U}_d$ is the uncertainty flag set.

\subsection{Required Data Variables and Units}
\label{sec:res-dat-unc-required-data-variables-units}

Mandatory fields are dataset/source name, provider, citation key, URL,
spatial coverage, temporal coverage, native resolution, horizontal CRS,
vertical datum, licence/usage note, project role, downstream
deliverable, preprocessing required and limitation.

\subsection{Candidate Data Sources}
\label{sec:res-dat-unc-candidate-data-sources}

% TODO: Complete this RES-DAT Domain-specific subsection.

\subsection{Spatial and Temporal Coverage}
\label{sec:res-dat-unc-spatial-temporal-coverage}

% TODO: Complete this RES-DAT Domain-specific subsection.

\subsection{Coordinate and Datum Requirements}
\label{sec:res-dat-unc-coordinate-datum-requirements}

No dataset shall be projected, merged or compared until its coordinate
and datum state is known or explicitly marked unresolved. Where a
source lacks confirmed CRS or datum metadata, the derived product shall
be flagged:

\begin{align}
  a_d
  &=
  \mathrm{screening\ only}
  \label{eq:res-dat-unc-screening-only-flag}
\end{align}

until the metadata are resolved.

\subsection{Data Processing and Transformation}
\label{sec:res-dat-unc-data-processing-transformation}

Preprocessing operations requiring explicit audit are reprojection,
vertical datum conversion, sign-convention conversion, tidal
correction, de-tiding, filtering, smoothing, gap filling, clipping,
raster/vector overlay, interpolation, mesh projection, source-to-bed
projection and replay-boundary sampling.

\subsection{Uncertainty, Provenance and Licensing}
\label{sec:res-dat-unc-uncertainty-provenance-licensing}

The uncertainty flag set shall include, where applicable:

\begin{itemize}
  \item source-age or survey-date heterogeneity;
  \item native-resolution mismatch across merged products;
  \item unknown or mixed vertical datums;
  \item horizontal CRS conversion uncertainty;
  \item smoothing or interpolation of sharp bathymetric/defence
    features;
  \item tide correction or harmonic-analysis dependence;
  \item source-inversion dependence for observation products;
  \item licence restrictions preventing redistribution;
  \item missing station, survey or defence metadata.
\end{itemize}

Source role: supports validation traceability and software ingestion.
Supporting dataset citations are recorded in `RES-DAT-CAS`.
Limitation: these flags identify uncertainty sources; they do not
replace quantitative uncertainty propagation.

\subsection{Data Acceptance Criteria}
\label{sec:res-dat-unc-data-acceptance-criteria}

A data asset is accepted for validation or solver input only when its
provenance record satisfies:

\begin{align}
  \mathcal{P}_{d}
  &\supseteq
  \left\{
    s_d,\,
    c_d,\,
    \mathcal{R}_{h,d},\,
    \mathcal{R}_{v,d},\,
    \Delta_d,\,
    \ell_d,\,
    \mathcal{T}_d
  \right\}
  \label{eq:res-dat-unc-minimum-accepted-provenance}
\end{align}

Records missing any of these fields may be used only for exploratory
screening unless the missing field is formally waived with a documented
project risk.

\subsection{Candidate Approaches}
\label{sec:res-dat-unc-candidate-approaches}

% TODO: Define the credible candidate approaches considered.

\subsection{Critical Comparison}
\label{sec:res-dat-unc-critical-comparison}

% TODO: Compare candidates using physically and project-relevant criteria.
% Do not select an approach solely on popularity or implementation ease.

\subsection{Assumptions and Applicability}
\label{sec:res-dat-unc-assumptions}

% TODO: State assumptions and the conditions under which they are valid.

\subsection{Limitations and Failure Conditions}
\label{sec:res-dat-unc-limitations}

% TODO: State where the evidence, model or selected approach ceases to be
% reliable or applicable.

\subsection{Project-Specific Conclusions}
\label{sec:res-dat-unc-conclusions}

The Tōhoku corridor comparison is sensitive to datum, resolution and
preprocessing differences. A single unqualified terrain raster, source
file or observation table is therefore not a valid project input. The
minimum accepted object is a data asset plus its provenance record
\(\mathcal{P}_{d}\).

\subsection{Selected Approach and Rationale}
\label{sec:res-dat-unc-selected-approach}

Use a manifest-first data workflow. Every source and derived product is
registered before use, and every transformation adds a new derived
record rather than overwriting the source identity.

\subsection{Unresolved Decisions}
\label{sec:res-dat-unc-unresolved-decisions}

Open decisions are final licence acceptance for J-EGG500/GSI-derived
products, exact vertical datum transformations, quantitative terrain
uncertainty values, observation station filtering rules and the
threshold for treating a missing metadata field as blocking rather than
screening-only.

\subsection{Requirements and Handoffs}
\label{sec:res-dat-unc-requirements}

Software shall implement a dataset manifest, immutable source records,
derived-product records, coordinate-transform audit logs, HDF5 or
equivalent structured output metadata and validation-report provenance
exports. `RES-VER-MET` shall consume the same uncertainty flags when
deciding whether a comparison is admissible.

\subsection{Deliverable Completion Record}
\label{sec:res-dat-unc-completion-record}

% TODO: Record whether the Deliverable is:
%
% - Not started
% - In progress
% - Draft complete
% - Reviewed
% - Accepted
%
% Acceptance requires:
%
% - the research questions to be answered;
% - claims to be supported by appropriate evidence;
% - assumptions and limitations to be explicit;
% - the project-specific conclusion to be stated;
% - downstream requirements to be traceable;
% - unresolved decisions to be closed or formally deferred.
